\chapter{Two-Dimensional Scattering Analyticity and Bootstrap Data}
\label{ch:integrable-scattering-analyticity-bootstrap}

\ConventionsLorentzian

The factorized scattering chapter fixed the algebraic data of an integrable
massive theory.  This chapter adds the analytic data in the rapidity plane.
The object studied here is an on-shell two-body scattering function together
with analytic hypotheses.  A local QFT realizing that function is a further
object; its construction requires Hilbert-space, locality, and domain data.

\section{Rapidity Plane and Physical Strip}

For equal masses \(m\), the two-particle Mandelstam invariant is
\[
  s(\theta)=2m^2(1+\cosh\theta)
  =4m^2\cosh^2\!\left(\frac{\theta}{2}\right).
\]
The physical two-particle scattering line is \(\theta\in\R\).  The crossed
channel is reached by \(\theta\mapsto \ii\pi-\theta\).  Hence the natural
analytic domain for a two-body amplitude is the physical strip
\[
  \mathcal S_{\rm phys}
  =
  \{\theta\in\C\mid 0<\operatorname{Im}\theta<\pi\}.
\]
For unequal masses \(m_a,m_b\),
\[
  s_{ab}(\theta)=m_a^2+m_b^2+2m_am_b\cosh\theta,
\]
and the same strip organizes the direct and crossed channels after the
species labels are crossed.

\begin{definition}[Factorized scattering bootstrap datum]
\label{def:factorized-scattering-bootstrap-datum}
A factorized scattering bootstrap datum is an elastic factorized scattering
datum of Definition~\ref{def:elastic-factorized-scattering-datum} together
with the following extra data and hypotheses:
\begin{enumerate}[label=(\roman*)]
\item meromorphic continuation of each component
  \(S_{ab}^{cd}(\theta)\) to the physical strip
  \(\mathcal S_{\rm phys}\);
\item a specified finite or locally finite list of simple-pole locations
  \(\theta=\ii u_{ab}^{e}\), \(0<u_{ab}^{e}<\pi\), labelled by proposed
  direct-channel bound-state species \(e\);
\item finite-dimensional multiplicity spaces for these bound-state channels
  and intertwiners \(\Gamma_{ab}^{e,\alpha}\) specifying residues;
\item horizontal-strip growth bounds stated as follows: for every closed
  substrip \(K\subset\mathcal S_{\rm phys}\) whose distance from the declared
  pole set is positive, there are constants \(C_K,N_K\) such that
  \[
    \|S_{ab}^{cd}(\theta)\|
    \le C_K(1+|\operatorname{Re}\theta|)^{N_K},
    \qquad \theta\in K.
  \]
  Near each declared simple pole the corresponding principal part is
  subtracted before applying this estimate.  When a thermodynamic
  Bethe-ansatz kernel
  \(\phi_{ab}(\theta)=-\ii\partial_\theta\log S_{ab}(\theta)\) is used later,
  the datum must also specify the branch and the additional integrability
  hypothesis on \(\phi_{ab}\), usually \(L^1(\mathbb R)\) after removal of
  explicitly treated CDD or bound-state factors.
\end{enumerate}
\end{definition}

The mass of a direct-channel bound state at \(\theta=\ii u_{ab}^{e}\) is
fixed by the pole location:
\[
  m_e^2
  =
  m_a^2+m_b^2+2m_am_b\cos u_{ab}^{e}.
\]
Indeed, write the incoming momenta as
\[
  p_a(\theta_a)=m_a(\cosh\theta_a,\sinh\theta_a),
  \qquad
  p_b(\theta_b)=m_b(\cosh\theta_b,\sinh\theta_b),
\]
with \(\theta=\theta_a-\theta_b\).  Then
\[
  (p_a+p_b)^2
  =
  m_a^2+m_b^2+2m_am_b\cosh(\theta_a-\theta_b)
  =
  m_a^2+m_b^2+2m_am_b\cosh\theta .
\]
Analytic continuation to the pole \(\theta=\ii u_{ab}^{e}\) gives
\(\cosh(\ii u_{ab}^{e})=\cos u_{ab}^{e}\), hence the displayed formula.  A
one-particle vector of mass \(m_e\) in a Hilbert-space QFT is additional data
to be constructed or assumed.  The pole and residue record the on-shell
criterion that such a vector would have to satisfy; a physical-strip pole
with the wrong residue sign, or a pole related by crossing rather than the
direct-channel residue, is not by itself a new stable direct-channel
particle.

\section{Residues and Fusion}

Let \(S_{ab}^{cd}(\theta)\) have a simple pole at
\(\theta=\ii u_{ab}^{e}\).  In a basis where the intermediate bound-state
space is explicitly included, the residue has the factorized form
\[
  -\ii\,\operatorname*{Res}_{\theta=\ii u_{ab}^{e}}
  S_{ab}^{cd}(\theta)
  =
  \sum_{\alpha}
  \Gamma_{ab}^{e,\alpha}\Gamma_{e,\alpha}^{cd}.
\]
The index \(\alpha\) labels a basis of the multiplicity space of species \(e\)
in the fusion of \(a\) and \(b\).  Positivity of the Hilbert-space norm fixes
the sign of the residue after the one-particle normalization has been chosen.

Fusion is the compatibility condition between a bound-state residue and
scattering with an additional particle \(k\).  If the rapidities of \(a\) and
\(b\) differ from the bound-state rapidity by imaginary fusing angles
\(\ii\bar u_{ae}^{b}\) and \(-\ii\bar u_{be}^{a}\), then the bootstrap
equation is the identity of meromorphic operator-valued functions
\[
  S_{ek}(\theta)
  =
  \Gamma_{ab}^{e}\,
  S_{ak}(\theta+\ii\bar u_{ae}^{b})\,
  S_{bk}(\theta-\ii\bar u_{be}^{a})\,
  \Gamma_{e}^{ab}.
\]
The products act on the tensor product of species spaces and multiplicity
spaces; the intertwiners \(\Gamma\) project onto the bound-state channel.  The
Yang--Baxter equation guarantees compatibility of this fused scattering with
three-particle factorization.

\begin{proposition}[Fusion identity as a residue identity]
\label{prop:fusion-identity-residue-derivation}
Assume that the three-particle scattering amplitude has meromorphic
factorization on the rapidity chamber under consideration and that the pair
\((a,b)\) has a simple direct-channel pole at
\(\theta_a-\theta_b=\ii u_{ab}^e\) with residue
\[
  -\ii\,\operatorname*{Res}_{\theta_a-\theta_b=\ii u_{ab}^{e}}
  S_{ab}(\theta_a-\theta_b)
  =
  \Gamma_{ab}^{e}\Gamma_{e}^{ab}.
\]
Choose the bound-state rapidity \(\theta\) so that the two constituent
rapidities are
\[
  \theta_a=\theta+\ii\bar u_{ae}^{b},
  \qquad
  \theta_b=\theta-\ii\bar u_{be}^{a},
  \qquad
  \bar u_{ae}^{b}+\bar u_{be}^{a}=u_{ab}^{e}.
\]
Then the residue of the factorized three-particle amplitude at the
\((a,b)\)-pole gives
\[
  S_{ek}(\theta-\theta_k)
  =
  \Gamma_{ab}^{e}\,
  S_{ak}(\theta-\theta_k+\ii\bar u_{ae}^{b})\,
  S_{bk}(\theta-\theta_k-\ii\bar u_{be}^{a})\,
  \Gamma_{e}^{ab}
\]
as an identity on the bound-state multiplicity space, after the
one-particle normalization used in the residue has been fixed.
\end{proposition}

\begin{proof}
Factorization writes the ordered three-particle scattering of
\((a,b,k)\) as a product in which the singular dependence on
\(\theta_a-\theta_b\) is carried by the two-body factor \(S_{ab}\), while
scattering with \(k\) is carried by \(S_{ak}\) and \(S_{bk}\) evaluated at the
displayed shifted rapidity differences.  Taking the residue at the
\((a,b)\)-pole therefore replaces the incoming pair by
\(\Gamma_{ab}^{e}\), the outgoing pair by \(\Gamma_e^{ab}\), and leaves the
two constituent scatterings with \(k\) between them.  The same residue, read
in the channel in which the pole has already been resolved as a one-particle
intermediate state \(e\), is by definition the scattering of \(e\) with
\(k\).  Equality of the two residues gives the formula.  The proof uses no
off-shell Lagrangian; all assumptions are meromorphic factorization, the
simple-pole residue normalization, and the existence of a one-particle
bound-state channel represented by the intertwiner \(\Gamma_{ab}^e\).
\end{proof}

\section{Minimal Solutions and CDD Factors}

For diagonal scattering, where \(S_{ab}^{cd}(\theta)=
\delta_a^c\delta_b^d S_{ab}(\theta)\), unitarity and crossing become
\[
  S_{ab}(\theta)S_{ab}(-\theta)=1,
  \qquad
  S_{ab}(\theta)=S_{\bar a b}(\ii\pi-\theta).
\]
A scalar meromorphic factor \(C(\theta)\) satisfying the same two equations
and having no assigned particle pole is a CDD factor.  Thus the analytic
constraints and the bound-state list determine the amplitude only after an
additional growth or minimality condition is specified.

\begin{definition}[Elementary scalar block]
\label{def:integrable-elementary-scalar-block}
For \(x\in\C\), the elementary scalar block is
\[
  [x]_\theta
  =
  \frac{\sinh\!\left((\theta+\ii\pi x)/2\right)}
       {\sinh\!\left((\theta-\ii\pi x)/2\right)} ,
\]
as a meromorphic function of \(\theta\).
\end{definition}

It satisfies \([x]_\theta[x]_{-\theta}=1\).  Products of such blocks place
zeros and poles at prescribed locations.  The bootstrap construction problem
is to choose the product and possible CDD factors so that unitarity, crossing,
physical-strip pole assignments, residue signs, and asymptotic bounds hold
simultaneously.

\begin{proposition}[Elementary block algebra]
\label{prop:elementary-scalar-block-algebra}
The elementary scalar block satisfies
\[
  [x]_{\theta}[x]_{-\theta}=1.
\]
Its poles are
\[
  \theta=\ii\pi x+2\pi\ii n,\qquad n\in\mathbb Z,
\]
and its zeros are
\[
  \theta=-\ii\pi x+2\pi\ii n,\qquad n\in\mathbb Z .
\]
For real \(0<x<1\), the pole at \(\theta=\ii\pi x\) lies in the physical
strip.
\end{proposition}

\begin{proof}
The unitarity identity is the cancellation
\[
  \frac{\sinh((\theta+\ii\pi x)/2)}
       {\sinh((\theta-\ii\pi x)/2)}
  \frac{\sinh((-\theta+\ii\pi x)/2)}
       {\sinh((-\theta-\ii\pi x)/2)}
  =
  \frac{\sinh((\theta+\ii\pi x)/2)}
       {\sinh((\theta-\ii\pi x)/2)}
  \frac{-\sinh((\theta-\ii\pi x)/2)}
       {-\sinh((\theta+\ii\pi x)/2)}
  =
  1.
\]
The denominator vanishes exactly when
\((\theta-\ii\pi x)/2=\pi\ii n\), and the numerator vanishes exactly when
\((\theta+\ii\pi x)/2=\pi\ii n\).  This gives the stated pole and zero sets.
\end{proof}

\section{Hilbert-Space Boundary}

The bootstrap datum is on-shell information.  A local QFT realization also
requires a Hilbert space, Wightman distributions, local operator domains, and
cluster properties.  In constructive treatments of factorizing models, wedge-
localized fields and modular localization provide one route from an
admissible \(S\)-matrix to local algebras.  The theorem boundary is precise:
analyticity and pole data enter as hypotheses for constructing wedge
operators and then proving that double-cone intersections are nontrivial.

\begin{openproblem}[Bootstrap completeness for local algebras]
\label{op:bootstrap-completeness-local-algebras}
Characterize the bootstrap data for which the resulting wedge-local
construction has nontrivial local algebras, a unique vacuum, asymptotic
completeness, and local fields whose form factors solve the equations of
Chapter~\ref{ch:integrable-factorized-scattering}.  This problem is the
operator-algebraic completion of the on-shell bootstrap.
\end{openproblem}
